\subsection*{Geofencing}
\label{sec:geofencing}

Geofencing polega na wykrywaniu przekroczenia przez obiekt wirtualnej granicy
geograficznej wyłącznie w oparciu o dane lokalizacyjne. Ostatecznie projekt ma na celu,
żeby użytkownik dostał powiadomienie w momencie, gdy węzeł przekroczy granicę wyznaczonego 
obszaru. W projekcie MeshTracker granice definiowane są jako wielokąty o dowolnym kształcie.

\subsection*{Model danych.}

Dane geofencingu przechowywane są w bazie Room za pomocą dwóch encji:
\texttt{Zone} i~\texttt{ZoneEvent} (tabele~\ref{tab:zone-entity}
i~\ref{tab:zoneevent-entity}).

\begin{table}[!h]
    \caption{Encja \texttt{Zone} - strefa geofencingu}
    \label{tab:zone-entity}
    \centering
    \begin{tabular}{|l|l|p{7cm}|}
        \hline
        \textbf{Pole} & \textbf{Typ} & \textbf{Opis} \\
        \hline
        \texttt{id}                 & \texttt{String (UUID)} & Klucz główny, generowany jako UUID v4 \\
        \texttt{name}               & \texttt{String}        & Nazwa nadana przez użytkownika \\
        \texttt{colorArgb}          & \texttt{Int}           & Kolor strefy (ARGB\footnote{ARGB - format kodowania koloru jako 32-bitowa liczba całkowita złożona z czterech 8-bitowych kanałów: Alpha (przezroczystość), Red, Green, Blue.} packed Int) \\
        \texttt{isActive}           & \texttt{Boolean}       & Czy strefa jest aktywnie monitorowana \\
        \texttt{verticesJson}       & \texttt{String}        & Wierzchołki wielokąta w formacie JSON \\
        \texttt{watchedNodeIdsJson} & \texttt{String}        & Lista ID węzłów do monitorowania w JSON \\
        \hline
    \end{tabular}
\end{table}

Wierzchołki wielokąta i lista monitorowanych węzłów przechowywane są jako serializowane
tablice JSON, przykładowo: \texttt{[\{"lat":52.1,"lon":21.0\},\ldots]}.

\begin{table}[!h]
    \caption{Encja \texttt{ZoneEvent} - zdarzenie wejścia/wyjścia ze strefy}
    \label{tab:zoneevent-entity}
    \centering
    \small
    \begin{tabular}{|p{2.8cm}|p{2.6cm}|p{6.5cm}|}
        \hline
        \textbf{Pole} & \textbf{Typ} & \textbf{Opis} \\
        \hline
        \texttt{id}               & \texttt{Long} (autoincrement) & Klucz główny \\
        \texttt{zoneId}           & \texttt{String}               & ID strefy  \\
        \texttt{nodeId}           & \texttt{String}               & ID węzła w chwili zdarzenia \\
        \texttt{eventType}        & \texttt{String}               & \texttt{"ENTER"} lub \texttt{"EXIT"} \\
        \texttt{\seqsplit{timestampSeconds}} & \texttt{Int}                  & Czas zdarzenia (Unix timestamp [s]) \\
        \hline
    \end{tabular}
    \normalsize
\end{table}


\paragraph{Algorytm sprawdzania przynależności do strefy.}

Do wykrywania, czy punkt GPS leży wewnątrz wielokąta strefy, zastosowano algorytm
\textbf{rzucania promienia} (ang.\ \textit{ray casting}) \footnote{\url{https://www.geeksforgeeks.org/cpp/point-in-polygon-in-cpp/}}.
Jest to jeden z~najszerzej stosowanych algorytmów rozwiązujących problem PIP (ang.\
\textit{Point-In-Polygon}) ze względu na prostotę implementacji i niską złożoność obliczeniową
$O(n)$, gdzie $n$ to liczba wierzchołków wielokąta~\cite{shimrat1962}. Uzasadnienie wyboru
akurat tej metody, oraz omówienie rozwiązań alternatywnych, zamieszczono na końcu rozdz.~\ref{sec:geofencing}.

Zasada działania jest następująca (rysunek~\ref{fig:ray-casting}). Z badanego punktu $P = (x_P, y_P)$
wysyłany jest poziomy promień w kierunku rosnącej długości geograficznej $y \to +\infty$.
Zliczana jest liczba przecięć promienia z krawędziami wielokąta:

\begin{itemize}
    \item nieparzysta liczba przecięć $\Rightarrow$ punkt leży \textbf{wewnątrz} wielokąta,
    \item parzysta liczba przecięć $\Rightarrow$ punkt leży \textbf{na zewnątrz}.
\end{itemize}

Dla każdej krawędzi łączącej wierzchołki $X_i = (x_i, y_i)$
i~$Y_j = (x_j, y_j)$ sprawdzany jest warunek przecięcia z promieniem:
krawędź musi mieć jeden wierzchołek ściśle powyżej szerokości geograficznej punktu
$x_P$, a drugi poniżej lub równo z nią. Jeśli warunek jest spełniony, wyznaczana jest
długość geograficzna przecięcia:

\begin{equation}
    y_{\text{przecięcia}} = y_i +
        \frac{x_P - x_i}{x_j - x_i} \cdot (y_j - y_i)
    \label{eq:ray-cast-cross}
\end{equation}

Przecięcie jest liczone tylko jeśli $y_P < y_{\text{przecięcia}}$, tzn.\ leży
na prawo od badanego punktu (promień biegnie w prawo). Każde takie przecięcie neguje
bieżącą wartość flagi \texttt{inside}, która na początku wynosi \texttt{false}.

\begin{figure}[!h]
    \label{fig:ray-casting}
    \centering
    \includegraphics[width=\linewidth]{strefaGeofencingRayCasting.jpg}
    \caption{Wizualizacja algorytmu ray casting - liczba punktów przecięcia z krawędziami jednoznacznie wskazuje, czy punkt leży wewnątrz czy na zewnątrz wielokąta}
\end{figure}


\begin{lstlisting}[caption={Implementacja algorytmu ray casting w \texttt{GeofenceChecker.contains()}},
                   label={lst:geofence-contains},
                   basicstyle=\footnotesize\ttfamily,
                   xleftmargin=0pt]
fun contains(
    pointLat: Double,
    pointLon: Double,
    polygon: List<ZoneVertex>
): Boolean {
    if (polygon.size < 3) return false

    var inside = false
    var j = polygon.lastIndex

    for (i in polygon.indices) {
        val vi = polygon[i]
        val vj = polygon[j]

        // Czy krawedz przecina poziomy promien z punktu w prawo?
        val crossesRay =
            (vi.lat <= pointLat && vj.lat > pointLat) ||
            (vj.lat <= pointLat && vi.lat > pointLat)

        if (crossesRay) {
            // Dlugosc geograficzna przeciecia promienia z krawedzia
            val intersectLon = vi.lon +
                (pointLat-vi.lat)/(vj.lat-vi.lat)*(vj.lon-vi.lon)

            if (pointLon < intersectLon) {
                inside = !inside
            }
        }
        j = i
    }
    return inside
}
\end{lstlisting}

Implementacja obsługuje wielokąty \emph{wklęsłe} i~\emph{wypukłe} oraz~\emph{złożone} (z przecięciami własnymi). 
Nie jest jednak zoptymalizowana pod kątem przypadków brzegowych, takich jak punkt leżący dokładnie na krawędzi 
lub wierzchołku wielokąta - w takich sytuacjach wynik może być nieprzewidywalny. W praktyce, ze względu na 
niedokładność GPS, takie przypadki są pomijalnie rzadkie i nie stanowią problemu dla funkcjonalności aplikacji.

\subsection*{Wykrywanie zdarzeń wejścia i wyjścia.}

Aby wygenerować zdarzenia \texttt{ENTER} i \texttt{EXIT}, konieczne jest śledzenie
zmian stanu między kolejnymi aktualizacjami pozycji.
\texttt{\seqsplit{ZoneMonitorService}} utrzymuje słownik \texttt{\seqsplit{previousNodeZoneMap}},
mapujący~ID~węzła na zbiór ID stref, w których węzeł znajdował się przy poprzedniej
aktualizacji. Po każdym odebraniu broadcastu \texttt{\seqsplit{ACTION\_NODE\_CHANGE}} wykonywane są trzy kroki:

\begin{enumerate}
    \item Obliczany jest zbiór stref \texttt{currentZones} zawierających węzeł
          w~chwili obecnej (tylko spośród stref, które ten węzeł monitorują).
    \item Zdarzenia \texttt{ENTER} generowane są dla stref obecnych w \texttt{currentZones},
          ale nieobecnych w \texttt{previousZones}.
    \item Zdarzenia \texttt{EXIT} generowane są dla stref obecnych w \texttt{previousZones},
          ale nieobecnych w \texttt{currentZones}.
\end{enumerate}

Po przetworzeniu słownik jest aktualizowany w taki sposób, że jeżeli węzeł nie znajduje się w~żadnej strefie,
jego wpis jest usuwany, co minimalizuje zużycie pamięci.

\paragraph{ZoneMonitorService - monitorowanie w tle.}

Logika geofencingu działa w ramach \texttt{ZoneMonitor\-Service} - usługi pierwszoplanowej
(ang.\ \textit{foreground service}) Androida. Jest to jedyna niezawodna metoda utrzymania
procesu aktywnego w tle: system Android agresywnie wstrzymuje zwykłe procesy w tle w celu
oszczędzania baterii, natomiast usługi pierwszoplanowe są przed tym chronione, pod warunkiem
że wyświetlają stałe powiadomienie systemowe~\cite{android-foreground-service}.

Cykl życia serwisu przebiega następująco:
\begin{itemize}
    \item \textbf{Uruchomienie}: \texttt{ZoneViewModel} wywołuje \texttt{ZoneMonitorService.start()},
          gdy w bazie danych istnieje \textbf{co najmniej jedna aktywna strefa}.
    \item \textbf{Działanie}: serwis rejestruje \texttt{MeshtasticBroadcastReceiver}
          i subskrybuje zmiany stref z bazy Room \footnote{Room to biblioteka Androida (część Jetpack) będąca warstwą abstrakcji nad SQLite. Zamiast pisać surowe zapytania SQL i ręcznie mapować wyniki na obiekty, definiuje się klasy encji (@Entity) i interfejsy DAO (@Dao) z adnotacjami, a Room generuje całą implementację w czasie kompilacji. W projekcie przechowuje strefy i log zdarzeń.} poprzez \texttt{Flow}\footnote{Flow to mechanizm z biblioteki Kotlin Coroutines reprezentujący asynchroniczny strumień wartości — analogia do LiveData, ale wbudowana w język. Room potrafi zwracać Flow<List<Zone>> zamiast jednorazowej listy: za każdym razem gdy dane w bazie się zmienią, Room automatycznie emituje nową listę do wszystkich subskrybentów. Właśnie dlatego zarówno ViewModel, jak i serwis są natychmiast powiadamiane o zmianach stref bez odpytywania bazy.}. Dzięki temu każda
          modyfikacja strefy jest propagowana do serwisu w czasie rzeczywistym, bez
          konieczności jego restartu.
    \item \textbf{Samozatrzymanie}: gdy \texttt{Flow} z Room zwróci pustą listę aktywnych
          stref (np.\ wszystkie zostały wyłączone lub usunięte), serwis wywołuje
          \texttt{stopSelf()} i kończy działanie autonomicznie.
    \item \textbf{Zatrzymanie zewnętrzne}: \texttt{ZoneViewModel} wywołuje
          \texttt{\seqsplit{ZoneMonitorService.stop()}} w reakcji na \textbf{wyłączenie ostatniej strefy}
          przez użytkownika.
\end{itemize}

\subsection*{Powiadomienia o przekroczeniu strefy.}

Wykrycie zdarzenia \texttt{ENTER} lub \texttt{EXIT} wyzwala wysłanie powiadomienia.
Powiadomienie zawiera nazwę węzła i nazwę strefy (rysunek~\ref{fig:strefa-powiadomienie-wejscia}), wysyłane jest z priorytetem \texttt{\seqsplit{PRIORITY\_HIGH}}
(kanał \texttt{\seqsplit{CHANNEL\_ZONE\_ALERTS}}) i~automatycznie znika po dotknięciu.
Przy pierwszym uruchomeniu aplikacji użytkownik musi jawnie przyznać
uprawnienie \texttt{\seqsplit{POST\_NOTIFICATIONS}}
(rysunek~\ref{fig:uprawnienia-powiadomienia}; obok lokalizacji i Bluetooth). Brak uprawnienia powoduje pominięcie powiadomienia bez
przerwania działania serwisu, co jest zachowaniem świadomym: serwis dalej wykrywa zdarzenia
i zapisuje je w bazie, jedynie nie wyświetla alertu.

\begin{figure}[!h]
    \centering
    \begin{minipage}{0.48\linewidth}
        \centering
        \includegraphics[width=\linewidth]{uprawnieniaPowiadomienia.jpg}
        \caption{Prośba o uprawnienie - \texttt{POST\_NOTIFICATIONS}}
        \label{fig:uprawnienia-powiadomienia}
    \end{minipage}
    \hfill
    \begin{minipage}{0.48\linewidth}
        \centering
        \includegraphics[width=\linewidth]{strefaPowiadomienieWejscia.jpg}
        \caption{Powiadomienie o wejściu w strefę}
        \label{fig:strefa-powiadomienie-wejscia}
    \end{minipage}
\end{figure}

Dotknięcie powiadomienia otwiera aplikację MeshTracker na ekranie mapy.

\paragraph{Uzasadnienie wyboru i metody alternatywne.}

Przed wyborem ray castingu rozważono kilka innych podejść, które warto
tu opisać - zarówno dlatego, że ich odrzucenie było nieoczywiste, jak i dlatego, że
w innych scenariuszach mogłyby być lepszym wyborem.

\textbf{Strefa okrągła (promień od centrum).}
Pierwszym pomysłem była strefa jako okrąg o środku $C = (x_C, y_C)$
i promieniu $r$. Sprawdzenie przynależności sprowadza się wtedy do prostego porównania:
\[
    d(P, C) \leq r \;\Leftrightarrow\; \text{punkt wewnątrz strefy}
\]
To rozwiązanie jest szybsze obliczeniowo i prostsze w implementacji, ponieważ nie trzeba
przechowywać listy wierzchołków. Rozwiązanie to odrzucono, ze względu na elastyczność. 
Zazwyczaj działki mają kształty zbliżone do wydłużonego wielokąta, a nie koła. 
Okrąg wystarczająco duży, żeby objąć cały teren, wychodziłby daleko poza drogę przy granicy działki, przez co alarm byłby
spóźniony. Okrąg wystarczająco mały, żeby nie wychodzić za drogę, nie obejmowałby
części terenu. Wielokąt pozwala narysować granicę dokładnie wzdłuż ogrodzenia sąsiada,
polnej miedzy, albo linii lasu.

\textbf{Google Geofencing API.}
Wstępnie rozważano skorzystanie z gotowego
geofencing API, które Android udostępnia właśnie do monitorowania przekroczeń granic
geograficznych. Szybko okazało się jednak, że to API monitoruje wyłącznie pozycję GPS
\emph{samego telefonu} i obsługuje tylko strefy okrągłe.


\textbf{Algorytm liczby owiniętej (ang.\ \textit{winding number}).}
W literaturze dotyczącej problemu PIP ~\cite{shimrat1962} zidentyfikowano również ten algorytm, który zlicza,
ile razy wielokąt owija się wokół badanego punktu. Jest numerycznie bardziej stabilny od
ray castingu, tzn.\ lepiej obsługuje przypadki brzegowe, gdzie punkt leży dokładnie na krawędzi
wielokąta. W kontekście niniejszego projektu rozróżnienie to jest jednak akademickie: Przy dokładności GPS rzędu kilku metrów prawdopodobieństwo, że zmierzony punkt
znajdzie się dokładnie na krawędzi wielokąta zdefiniowanego przez użytkownika, jest
pomijalnie małe. Wybór padł na ray casting, ponieważ implementacja jest krótsza, prostsza do zrozumienia
i wystarczająca na potrzeby projektu. Złożoność obliczeniowa obu algorytmów wynosi $O(n)$.

\textbf{Indeksowanie przestrzenne (R-tree\footnote{R-tree - hierarchiczna struktura danych do indeksowania obiektów przestrzennych, grupująca je w zagnieżdżone prostokąty otaczające. Umożliwia szybkie wyszukiwanie przestrzenne w czasie $O(\log n)$.}).}
W przypadku konieczności obsługiwania setek stref jednocześnie, sprawdzanie każdej z nich przy 
aktualizacji pozycji stanowiłoby istotne ograniczenie wydajności.
Struktury takie jak R-tree pozwalają szybko zawęzić kandydatów do stref, których obwiednia
prostokątna zawiera dany punkt. W MeshTrackerze spodziewanych jest co najwyżej kilka stref
i kilka węzłów - indeksowanie nie przyniosłoby żadnych mierzalnych korzyści, a znacznie
skomplikowałoby implementację.
